\newpage{}
\section{From Edges to Nodes}
\label{sec:edges-to-nodes}

This section defines a concrete procedure for lifting edge-carried market
features into node-carried market features.  We call this procedure
\[
  \operatorname{NodeLift}.
\]
It is not a forecast model.  It is a graph-structured transformation applied
to observed input windows.

The guiding idea is:

\begin{quote}
Edges carry market relations and market activity.  Nodes carry recovered
semantic state.  A node lift turns edge-local observations into node-local
features, while preserving residual edge information that cannot be explained by
nodes alone.
\end{quote}

For one anchor time $t$, one channel $c$, one input position $i$, and one
directed edge $e=(u,v)$, write the observed edge normalized feature vector as
\[
  \tensor{x}_{e}^{c\,i}(t)
  :=
  \tensor{x}_{u\,v}^{c\,i}(t)
  \in
  \R^{D_x}.
\]
The lift produces a node feature matrix
\[
  \tensor{s}^{c\,i}(t)
  \in
  \R^{N\times D_y},
\]
where $s_u^{c\,i\,d}(t)$ is the lifted feature coordinate attached to node
$u\in\mathcal{V}$.

Here $D_x$ is the source kline feature width before lifting, and $D_y$ is the
node feature width after lifting.  The symbol $D_y$ is kept because it names the
output side of NodeLift.  For the full normalized kline lift, one node feature
is produced for each input feature, so
\[
  D_y=D_x=9.
\]
If a later model selects only a subset of node coordinates, that selection
happens after this input-side lift.

\calloutbox{NodeLift principle.}{\\%
  Price coordinates are relative quantities, so they are lifted through signed
  edge differences.  Activity coordinates are additive quantities, so they are
  lifted through endpoint aggregation.  The residual edge field is preserved
  because it contains cycle tension and edge-local information not represented
  by the node state.
}

\subsection{Feature Classes}
\label{subsec:node-lift-feature-classes}

The normalized kline registry has four price-return coordinates and five
activity coordinates.  We write
\[
  \mathcal{J}_{\mathrm{price}}
  =
  \{1,2,3,4\},
\]
corresponding to \field{open}, \field{high}, \field{low}, and
\field{close}, and
\[
  \mathcal{J}_{\mathrm{act}}
  =
  \{5,6,7,8,9\},
\]
corresponding to \field{volume}, \field{quote\_volume}, \field{trades},
\field{taker\_buy\_base}, and \field{taker\_buy\_quote}.

The price coordinates are normalized log-return coordinates.  They are
compatible with a difference equation on node log-valuations.  The activity
coordinates are normalized nonnegative activity coordinates.  They are
compatible with endpoint aggregation, not with a difference equation.

Time keys such as \field{open\_time} and \field{close\_time} are alignment
metadata rather than model features.  They are handled by the sequence
synchronization procedure and are not lifted here.

The normalization rules used by these coordinates are recorded next, because
NodeLift depends on the difference between price-like coordinates and
activity-like coordinates.

\subsection{Normalization Procedure Used by NodeLift}
\label{subsec:nodelift-normalization-procedure}

NodeLift operates on normalized edge features.  It does not lift raw kline
records directly.  For one directed edge $e=(u,v)$, channel $c$, and input
position $i$, the normalized slice is
\[
  \tensor{x}_{e}^{c\,i}(t)
  \in
  \R^{D_x},
  \qquad
  D_x=9.
\]
The corresponding input mask is
\[
  {}_{x}m_{e}^{c\,i}(t)
  \in
  \{0,1\}.
\]
A masked-out slice is not a zero-valued market event.  It is unavailable
evidence.

The normalization pass is causal and channel-local.  For each edge $e$ and
channel $c$, the normalizer scans records from past to future while remembering
the previous valid candle on that same edge and channel.  Let
\[
  p_{e}^{c\,(i-1)\,(\mathrm{close})}(t)\in\Rpos
\]
denote the previous valid close price available before position $i$ on the same
edge and channel.

For one raw valid candle, write the OHLC price tuple for $\mathcal{J}_{\mathrm{price}}$ as
\[
  \left(
    p_{e}^{c\,i\,(\mathrm{open})},
    p_{e}^{c\,i\,(\mathrm{high})},
    p_{e}^{c\,i\,(\mathrm{low})},
    p_{e}^{c\,i\,(\mathrm{close})}
  \right)
  \in
  \Rpos^4,
\]
and the raw activity tuple for $\mathcal{J}_{\mathrm{act}}$ as
\[
  \left(
    a_{e}^{c\,i\,(\mathrm{volume})},
    a_{e}^{c\,i\,(\mathrm{quote})},
    a_{e}^{c\,i\,(\mathrm{trades})},
    a_{e}^{c\,i\,(\mathrm{buybase})},
    a_{e}^{c\,i\,(\mathrm{buyquote})}
  \right)
  \in
  \Rnonneg^5.
\]

If the current raw candle is invalid, the normalized slice is marked invalid:
\[
  {}_{x}m_{e}^{c\,i}(t)=0.
\]
The previous-valid state is cleared, so the next valid candle begins a new
contiguous valid segment.

If the current raw candle is valid but no previous valid close is available,
the normalized slice is also marked invalid:
\[
  {}_{x}m_{e}^{c\,i}(t)=0.
\]
This is why the first valid candle of every contiguous valid segment is masked
out: there is no causal previous close from which to form log-return price
coordinates.  The current valid candle may then become the previous valid candle
for the next position.

If both the current raw candle and the previous valid close are available, the
normalized price coordinates $d\in\mathcal{J}_{\mathrm{price}}$ are
\[
  x_{e}^{c\,i\,(1)}(t)
  =
  \log
  \left(
    \frac{
      p_{e}^{c\,i\,(\mathrm{open})}
    }{
      p_{e}^{c\,(i-1)\,(\mathrm{close})}(t)
    }
  \right),
\]
\[
  x_{e}^{c\,i\,(2)}(t)
  =
  \log
  \left(
    \frac{
      p_{e}^{c\,i\,(\mathrm{high})}
    }{
      p_{e}^{c\,(i-1)\,(\mathrm{close})}(t)
    }
  \right),
\]
\[
  x_{e}^{c\,i\,(3)}(t)
  =
  \log
  \left(
    \frac{
      p_{e}^{c\,i\,(\mathrm{low})}
    }{
      p_{e}^{c\,(i-1)\,(\mathrm{close})}(t)
    }
  \right),
\]
and
\[
  x_{e}^{c\,i\,(4)}(t)
  =
  \log
  \left(
    \frac{
      p_{e}^{c\,i\,(\mathrm{close})}
    }{
      p_{e}^{c\,(i-1)\,(\mathrm{close})}(t)
    }
  \right).
\]

The normalized activity coordinates $d\in\mathcal{J}_{\mathrm{act}}$ are
\[
  x_{e}^{c\,i\,(5)}(t)
  =
  \log
  \left(
    1+a_{e}^{c\,i\,(\mathrm{volume})}
  \right),
\]
\[
  x_{e}^{c\,i\,(6)}(t)
  =
  \log
  \left(
    1+a_{e}^{c\,i\,(\mathrm{quote})}
  \right),
\]
\[
  x_{e}^{c\,i\,(7)}(t)
  =
  \log
  \left(
    1+a_{e}^{c\,i\,(\mathrm{trades})}
  \right),
\]
\[
  x_{e}^{c\,i\,(8)}(t)
  =
  \log
  \left(
    1+a_{e}^{c\,i\,(\mathrm{buybase})}
  \right),
\]
and
\[
  x_{e}^{c\,i\,(9)}(t)
  =
  \log
  \left(
    1+a_{e}^{c\,i\,(\mathrm{buyquote})}
  \right).
\]

When all nine normalized coordinates are produced successfully, the slice is
valid:
\[
  {}_{x}m_{e}^{c\,i}(t)=1.
\]
The resulting normalized vector is
\[
  \tensor{x}_{e}^{c\,i}(t)
  =
  \left(
    x_{e}^{c\,i\,(1)}(t),
    \ldots,
    x_{e}^{c\,i\,(9)}(t)
  \right)
  \in
  \R^9.
\]

In implementation, the normalization pass protects the domains of the
logarithms.  Nonpositive or non-finite prices are clamped to a positive floor
before price logarithms are applied, and non-finite activity values are clamped
to a valid nonnegative value before $\log(1+\cdot)$ is applied.  These numerical
protections keep the normalized cache finite, but they do not turn invalid
records into observed market evidence.  The mask remains the source of truth
for whether a normalized slice is usable.

\calloutbox{Normalization and NodeLift.}{\\%
  NodeLift reads the normalized feature coordinates.  The price coordinates
  $d\in\mathcal{J}_{\mathrm{price}}$ are log-return coordinates and are lifted
  by signed incidence recovery.  The activity coordinates
  $d\in\mathcal{J}_{\mathrm{act}}$ are $\log(1+a)$ coordinates and are lifted
  by first applying the inverse activity map $a=\exp(x)-1$, then aggregating to
  endpoints.
}

\subsection{Balance Equation for Price Coordinates}
\label{subsec:balance-equation-for-price-coordinates}

The price-coordinate lift comes from the usual edge-to-node balance relation.
Suppose every node $u\in\mathcal{V}$ has a positive latent valuation
\[
  \eta_u\in\Rpos.
\]
For an edge $e=(u,v)$, the idealized edge price is the ratio
\[
  q_e
  =
  \frac{\eta_u}{\eta_v}
  \in\Rpos.
\]
In log coordinates,
\[
  \xi_u := \log\eta_u\in\R,
  \qquad
  \ell_e := \log q_e\in\R,
\]
so
\[
  \ell_e
  =
  \xi_u-\xi_v.
\]
If the node log-valuations change over time, define the node increment
\[
  y_u := \Delta\xi_u\in\R.
\]
Then the scalar edge return satisfies
\[
  r_e
  =
  \Delta\ell_e
  =
  y_u-y_v.
\]

This is the balance equation:
\[
  r_e = y_u-y_v.
\]
The edge return is a local relation.  The node vector is the global potential
whose differences explain those local relations.

\subsection{Observed Input Masks and Usable Edge Sets}
\label{subsec:observed-input-masks-and-usable-edge-sets}

NodeLift must lift masks together with values.  The input mask
\[
  {}_{x}m_{e}^{c\,i}(t)
  \in
  \{0,1\},
  \qquad
  e=(u,v),
\]
says whether the normalized edge feature slice
$\tensor{x}_{e}^{c\,i}(t)\in\R^{D_x}$ is valid.  A masked-out slice is not an
observed zero-valued market event.  It is unavailable evidence.

Since NodeLift produces node features with a feature coordinate, we first
derive a coordinatewise edge mask.  By default,
\[
  {}_{x}m_{e}^{c\,i\,d}(t)
  :=
  {}_{x}m_{e}^{c\,i}(t)
  \cdot
  \mathbf{1}
  \left\{
    x_{e}^{c\,i\,d}(t)
    \text{ is valid}
  \right\}
  \in
  \{0,1\}.
\]
In the normalized cache this valid-value check is usually already protected by
the normalization procedure, because the cache stores finite normalized values,
but it is useful to keep the check in the mathematical contract.  Later
feature-level filters may also refine ${}_{x}m_{e}^{c\,i\,d}(t)$ without
changing the rest of the lift.

The slice-level usable edge set is
\[
  \mathcal{E}_{\mathrm{slice}}^{c\,i}(t)
  =
  \left\{
    e\in\mathcal{E}
    :
    {}_{x}m_{e}^{c\,i}(t)=1
  \right\}.
\]
For each coordinate, define the coordinatewise usable edge set
\[
  \mathcal{E}_{\mathrm{use}}^{c\,i\,d}(t)
  =
  \left\{
    e\in\mathcal{E}
    :
    {}_{x}m_{e}^{c\,i\,d}(t)=1
  \right\}
  \subseteq
  \mathcal{E}_{\mathrm{slice}}^{c\,i}(t).
\]
Thus usability is coordinatewise.  The usable edge set for a price coordinate
and the usable edge set for an activity coordinate need not be identical.

For observed input windows, the edge-coordinate field over the coordinatewise
usable edge set is
\[
  \vect{x}_{\mathcal{E}}^{c\,i\,d}(t)
  =
  \left(
    x_{e}^{c\,i\,d}(t)
  \right)_{e\in\mathcal{E}_{\mathrm{use}}^{c\,i\,d}(t)}
  \in
  \R^{L_{\mathrm{use}}^{c\,i\,d}},
\]
where
\[
  L_{\mathrm{use}}^{c\,i\,d}
  =
  \left|
    \mathcal{E}_{\mathrm{use}}^{c\,i\,d}(t)
  \right|.
\]
This is the scalar field consumed by NodeLift in this chapter:
\[
  \text{observed scalar input edge field}
  \quad
  \longrightarrow
  \quad
  \text{node field}.
\]
Predictive and supervised readouts may reuse the same algebra later, but they
should introduce their source explicitly instead of hiding it behind a generic
edge-field symbol.

\subsection{Graph Matrices for NodeLift}
\label{subsec:graph-matrices-for-nodelift}

The graph matrices used by NodeLift are built from the endpoints of each
directed edge.  For every edge
\[
  e=(u(e),v(e)),
  \quad
  e\in\mathcal{E}.
\]
Where $u(e)\in\mathcal{V}$ is the base endpoint and
$v(e)\in\mathcal{V}$ is the quote endpoint.
Thus, for an edge such as \field{BTCUSDT},
$u(e)=\text{\field{BTC}}$, $v(e)=\text{\field{USDT}}$.

The graph matrices have one row per edge and one column per node.  In this
subsection, the symbol $a\in\mathcal{V}$ denotes a generic node column.  It is
not a raw activity variable.

The signed incidence matrix of the active graph is
\[
  \mat{A}_{\mathcal{E}}
  \in
  \R^{L\times N},
\]
with entries%
\note{
  The symbol $a\in\mathcal{V}$ denotes a generic node column. It is not a raw activity variable.
}
\[
  A_{e}^{a}
  =
  \mathbf{1}\{a=u(e)\}
  -
  \mathbf{1}\{a=v(e)\},
  \qquad
  e\in\mathcal{E},
  \quad
  a\in\mathcal{V}.
\]
Therefore, for any node vector $\vect{y}\in\R^N$,
\[
  \left(
    \mat{A}_{\mathcal{E}}\vect{y}
  \right)_{e}
  =
  \sum_{a\in\mathcal{V}}
  A_{e}^{a}y_a
  =
  y_{u(e)}-y_{v(e)}.
\]
This is the signed operator used for price-like coordinates.

For activity aggregation, define the base-endpoint assignment matrix
\[
  \mat{B}_{\mathcal{E}}
  \in
  \{0,1\}^{L\times N}
\]
by
\[
  B_{e}^{a}
  =
  \mathbf{1}\{a=u(e)\},
  \qquad
  e\in\mathcal{E},
  \quad
  a\in\mathcal{V}.
\]
Thus $B_{e}^{a}=1$ exactly when node $a$ is the base endpoint of edge $e$.

Define the quote-endpoint assignment matrix
\[
  \mat{Q}_{\mathcal{E}}
  \in
  \{0,1\}^{L\times N}
\]
by
\[
  Q_{e}^{a}
  =
  \mathbf{1}\{a=v(e)\},
  \qquad
  e\in\mathcal{E},
  \quad
  a\in\mathcal{V}.
\]
Thus $Q_{e}^{a}=1$ exactly when node $a$ is the quote endpoint of edge $e$.

The unsigned endpoint matrix is
\[
  \mat{U}_{\mathcal{E}}
  =
  \mat{B}_{\mathcal{E}}
  +
  \mat{Q}_{\mathcal{E}}
  \in
  \{0,1\}^{L\times N}.
\]
Its entries are
\[
  U_{e}^{a}
  =
  \mathbf{1}\{a=u(e)\}
  +
  \mathbf{1}\{a=v(e)\}.
\]
Since self-edges are not allowed, $u(e)\ne v(e)$, so
$U_{e}^{a}$ is either $0$ or $1$.

These matrices have simple aggregation meanings.  For any scalar edge field
\[
  \vect{\rho}_{\mathcal{E}}
  =
  (\rho_e)_{e\in\mathcal{E}}
  \in
  \R^{L},
\]
the base aggregation is
\[
  \left(
    \mat{B}_{\mathcal{E}}^{\top}
    \vect{\rho}_{\mathcal{E}}
  \right)^{a}
  =
  \sum_{e\in\mathcal{E}}
  B_{e}^{a}\rho_e
  =
  \sum_{\substack{e\in\mathcal{E}\\u(e)=a}}
  \rho_e.
\]
The quote aggregation is
\[
  \left(
    \mat{Q}_{\mathcal{E}}^{\top}
    \vect{\rho}_{\mathcal{E}}
  \right)^{a}
  =
  \sum_{e\in\mathcal{E}}
  Q_{e}^{a}\rho_e
  =
  \sum_{\substack{e\in\mathcal{E}\\v(e)=a}}
  \rho_e.
\]
The incident endpoint aggregation is
\[
  \left(
    \mat{U}_{\mathcal{E}}^{\top}
    \vect{\rho}_{\mathcal{E}}
  \right)^{a}
  =
  \sum_{\substack{e\in\mathcal{E}\\u(e)=a\ \mathrm{or}\ v(e)=a}}
  \rho_e.
\]

Thus:
\[
  \mat{A}_{\mathcal{E}}
  \quad
  \text{creates signed edge differences from node values},
\]
while
\[
  \mat{B}_{\mathcal{E}},
  \quad
  \mat{Q}_{\mathcal{E}},
  \quad
  \mat{U}_{\mathcal{E}}
  \quad
  \text{assign edge activity to node endpoints}.
\]

For a coordinatewise usable edge set
\[
  \mathcal{E}_{\mathrm{use}}^{c\,i\,d}(t)
  \subseteq
  \mathcal{E},
\]
the corresponding usable graph matrices are obtained by keeping only the rows
whose edge belongs to that set:
\[
  \mat{A}^{c\,i\,d}(t)
  =
  \left.
    \mat{A}_{\mathcal{E}}
  \right|_{\mathcal{E}_{\mathrm{use}}^{c\,i\,d}(t)},
\]
\[
  \mat{B}^{c\,i\,d}(t)
  =
  \left.
    \mat{B}_{\mathcal{E}}
  \right|_{\mathcal{E}_{\mathrm{use}}^{c\,i\,d}(t)},
  \qquad
  \mat{Q}^{c\,i\,d}(t)
  =
  \left.
    \mat{Q}_{\mathcal{E}}
  \right|_{\mathcal{E}_{\mathrm{use}}^{c\,i\,d}(t)},
\]
and
\[
  \mat{U}^{c\,i\,d}(t)
  =
  \left.
    \mat{U}_{\mathcal{E}}
  \right|_{\mathcal{E}_{\mathrm{use}}^{c\,i\,d}(t)}.
\]
Here the vertical restriction symbol means: keep the selected edge rows and keep
all node columns.

Once the channel $c$, input position $i$, coordinate $d$, and usable edge set
are fixed, we suppress these labels and write simply
\[
  \mat{A},
  \qquad
  \mat{B},
  \qquad
  \mat{Q},
  \qquad
  \mat{U}.
\]

\calloutbox{Endpoint matrices.}{\\%
  The matrix $\mat{A}$ is signed: it compares the base endpoint against the
  quote endpoint.  The matrices $\mat{B}$ and $\mat{Q}$ are unsigned assignment
  matrices: $\mat{B}$ sends an edge quantity to its base node, and $\mat{Q}$
  sends an edge quantity to its quote node.  The matrix $\mat{U}=\mat{B}+\mat{Q}$
  sends an edge quantity to both endpoints.
}

\subsection{Edge Conditions for a Valid Lift}
\label{subsec:edge-conditions-for-valid-lift}

NodeLift does not require the active graph to be complete.  The active edge set
may be any staged set of traded instruments:
\[
  \mathcal{E}
  \subseteq
  \{(u,v)\in\mathcal{V}\times\mathcal{V}:u\ne v\}.
\]
For one channel, position, and coordinate, the usable edge set
\[
  \mathcal{E}_{\mathrm{use}}^{c\,i\,d}(t)
  \subseteq
  \mathcal{E}
\]
may be smaller after masks, data availability, feature availability, and
observer choices are applied.

For price-coordinate recovery, the important condition is connectivity of the
underlying undirected usable graph.  Direction fixes the sign of each incidence
row, but the ability to recover node differences depends on whether the usable
edges connect the nodes.

A reverse edge is not required.  If $e=(u,v)$ is present, then its incidence row
already gives the relation
\[
  y_u-y_v.
\]
If the reverse edge $(v,u)$ is also present, it is a second directed
observation with the opposite incidence sign:
\[
  y_v-y_u.
\]
Both may be used when both are independently observed traded instruments.  But
if one reverse edge is synthetically created from the other, it should not be
treated as an independent observation with full confidence, because that would
double-count the same market information.

Let the usable undirected graph for one price coordinate have connected
components
\[
  \mathcal{V}_{(1)},\ldots,\mathcal{V}_{(C_{\mathrm{comp}})}.
\]
For component $a$, write
\[
  n_a = |\mathcal{V}_{(a)}|,
  \qquad
  L_a = |\mathcal{E}_{(a)}|.
\]
For price recovery on that component, the component must contain enough usable
edges to connect its nodes.  At minimum, it needs
\[
  L_a \ge n_a-1
\]
with those edges arranged as a connected tree.  If fewer edges are available, the
price-coordinate node values inside that component are underdetermined and the
missing nodes or component should be masked out.

A connected tree is enough to recover node coordinates after a gauge is fixed.
However, a tree has no cycles.  Therefore it cannot reveal cycle inconsistency:
any edge field on a tree can be explained exactly by some node-potential vector.
Extra edges beyond a tree create cycles.  Cycles make the edge field
overdetermined and allow residual tension to appear.

For one connected component, the number of independent cycle constraints is
\[
  L_a-n_a+1.
\]
Across all usable components, the total cycle dimension is
\[
  L_{\mathrm{use}}
  -
  N_{\mathrm{use}}
  +
  C_{\mathrm{comp}},
\]
where $N_{\mathrm{use}}$ is the number of nodes touched by the usable edge set.
When this number is zero, the usable graph is a forest.  A forest can still
recover node coordinates componentwise, but it cannot detect cycle
inconsistency.  Nonzero residual tension requires cycles.

Activity-coordinate lifting has weaker graph requirements.  A node can receive
an activity feature whenever the relevant endpoint aggregation has usable
incident edges.  Activity aggregation does not require a connected price graph.

\calloutbox{Edge requirement.}{\\%
  NodeLift does not need all possible instruments and does not need both
  directions of every pair.  For price coordinates, each recoverable component
  needs enough usable edges to be connected.  Extra edges beyond a tree create
  cycles, and cycles are where residual tension and arbitrage-like disagreement
  can appear.
}

\subsection{Synthetic Reference Gauge}
\label{subsec:synthetic-reference-gauge}

Because only node differences appear, adding the same constant to every node in
one connected component does not change any edge difference:
\[
  (y_u-y_v)
  =
  \bigl((y_u+a)-(y_v+a)\bigr),
  \qquad
  a\in\R.
\]
So price-coordinate recovery is not unique until a reference is fixed.

The ordinary gauge fixes one real node $g\in\mathcal{V}$ by imposing
\[
  y_g=0.
\]
That is valid, but it makes every displayed node coordinate relative to that
single chosen node.

Instead, NodeLift uses a synthetic reference vector.  For one connected usable
component, let
\[
  \vect{w}_{\star}
  \in
  \Rnonneg^{n_a},
  \qquad
  \mathbf{1}^{\top}\vect{w}_{\star}=1.
\]
The synthetic reference return is
\[
  y_{\star}
  =
  \vect{w}_{\star}^{\top}\vect{y},
\]
and the gauge condition is
\[
  \vect{w}_{\star}^{\top}\vect{y}=0.
\]

A real-node gauge is a special case.  If $\vect{w}_{\star}$ places all mass on
node $g$, then
\[
  \vect{w}_{\star}^{\top}\vect{y}
  =
  y_g.
\]
The synthetic gauge therefore generalizes the ordinary constraint $y_g=0$.

For multiple connected components, one independent reference is needed per
component.  Let
\[
  \vect{w}_{\star}^{(a)}
\]
be the synthetic reference for component $a$, supported only on
$\mathcal{V}_{(a)}$.  Stack these component references into the reference
constraint matrix
\[
  \mat{G}_{\star}
  =
  \begin{bmatrix}
    \vect{w}_{\star}^{(1)}
    &
    \cdots
    &
    \vect{w}_{\star}^{(C_{\mathrm{comp}})}
  \end{bmatrix}
  \in
  \R^{N\times C_{\mathrm{comp}}}.
\]
The multi-component gauge condition is
\[
  \mat{G}_{\star}^{\top}\vect{y}
  =
  \vect{0}.
\]
For a connected usable graph, this reduces to the single constraint
\[
  \vect{w}_{\star}^{\top}\vect{y}=0.
\]

In log-valuation coordinates, the synthetic node is a geometric basket:
\[
  \xi_{\star}
  =
  \sum_{u\in\mathcal{V}}
  w_{\star,u}\xi_u,
  \qquad
  \eta_{\star}
  =
  \prod_{u\in\mathcal{V}}
  \eta_u^{w_{\star,u}}.
\]
Thus the synthetic reference is not a stablecoin and does not need to be a
traded asset.  It is the coordinate origin from which recovered node values are
read.

\calloutbox{Gauge interpretation.}{\\%
  The synthetic reference chooses the coordinate origin for recovered
  price-coordinate node values.  It can shift the displayed price-coordinate
  values within a connected component, but it does not change edge differences
  and does not erase residual edge tension.  Activity coordinates are endpoint
  aggregates rather than edge differences, so this gauge does not act on them.
  Residuals remain the gauge-invariant place where cycle disagreement appears.
}

\subsection{Choosing the Covariance-Minimizing Reference}
\label{subsec:choosing-covariance-minimizing-reference}

When a node covariance estimate is available, the synthetic reference may be
chosen to make displayed node coordinates as stable as possible.

For one connected component, define the admissible set of reference baskets
\[
  \mathcal{W}_{\star}
  =
  \left\{
    \vect{w}\in\Rnonneg^{n_a}:
    \mathbf{1}^{\top}\vect{w}=1
  \right\}.
\]
Let
\[
  \mat{C}_{y}^{\mathrm{quot}}
\]
denote a quotient-space node covariance for that component, where the additive
constant direction has already been identified as a gauge symmetry.  This
covariance may come from a previous recovery pass, repeated observations,
sampled scenarios, predictive ensembles, or a regularized graph covariance
estimator.

For any candidate reference $\vect{w}\in\mathcal{W}_{\star}$, define the
centering map
\[
  \mat{P}_{w}
  :=
  \mat{I}
  -
  \mathbf{1}\vect{w}^{\top}.
\]
It transforms a node vector into coordinates read relative to $\vect{w}$:
\[
  \widetilde{\vect{y}}
  =
  \mat{P}_{w}\vect{y}
  =
  \vect{y}
  -
  \mathbf{1}
  \left(
    \vect{w}^{\top}\vect{y}
  \right),
  \qquad
  \vect{w}^{\top}\widetilde{\vect{y}}=0.
\]
The displayed covariance under this reference is
\[
  \mat{C}_{y}^{(w)}
  =
  \mat{P}_{w}
  \mat{C}_{y}^{\mathrm{quot}}
  \mat{P}_{w}^{\top}.
\]

Let $\mat{D}_{\mathrm{node}}$ be an optional diagonal matrix of node-importance
weights.  If no priority is needed, take
\[
  \mat{D}_{\mathrm{node}}=\mat{I}.
\]
Let $\vect{w}_0\in\mathcal{W}_{\star}$ be an optional prior reference basket,
often initialized as
\[
  \vect{w}_0
  =
  \frac{1}{n_a}\mathbf{1}.
\]
The covariance-minimizing reference is
\[
  \vect{w}_{\star}
  =
  \operatorname*{arg\,min}_{\vect{w}\in\mathcal{W}_{\star}}
  \left\{
    \operatorname{tr}
    \left(
      \mat{D}_{\mathrm{node}}
      \mat{P}_{w}
      \mat{C}_{y}^{\mathrm{quot}}
      \mat{P}_{w}^{\top}
    \right)
    +
    \lambda_{\mathrm{ref}}
    \left\|
      \vect{w}-\vect{w}_0
    \right\|_2^2
  \right\},
  \qquad
  \lambda_{\mathrm{ref}}\in\Rnonneg.
\]
This chooses the synthetic reference whose displayed node coordinates have small
covariance, while optionally staying close to a prior basket.

If no covariance estimate is available yet, one may initialize
\[
  \vect{w}_{\star}
  =
  \frac{1}{n_a}\mathbf{1}
\]
on each connected component, or use a selected real node as a temporary
reference.  The recovery equations below only require valid reference
constraints, one per usable connected component.

\subsection{Price-Coordinate Lift}
\label{subsec:price-coordinate-lift}

For every price coordinate
\[
  d\in\mathcal{J}_{\mathrm{price}},
\]
the edge coordinate is interpreted as a relative log-price feature.  The node
lift asks for node values whose differences explain the edge field:
\[
  \vect{x}_{\mathcal{E}}^{c\,i\,d}
  \approx
  \mat{A}\vect{y}^{c\,i\,d}.
\]

Let
\[
  \mat{W}^{c\,i\,d}
  \in
  \R^{L_{\mathrm{use}}^{c\,i\,d}\times L_{\mathrm{use}}^{c\,i\,d}}
\]
be the edge precision matrix for this coordinate.  In the simplest independent
case,
\[
  \mat{W}^{c\,i\,d}
  =
  \mat{\Gamma}^{c\,i\,d}
  =
  \diag
  \left(
    \gamma_1^{c\,i\,d},
    \ldots,
    \gamma_{L_{\mathrm{use}}}^{c\,i\,d}
  \right).
\]
If a full edge covariance estimate is available, then
\[
  \mat{W}^{c\,i\,d}
  =
  \left(
    \mat{C}_{x}^{c\,i\,d}
  \right)^{-1},
\]
using a regularized inverse or pseudoinverse when needed.

The recovered node price-coordinate vector is
\[
  \widehat{\vect{y}}^{c\,i\,d}
  =
  \operatorname*{arg\,min}_{\vect{y}\in\R^N}
  \left(
    \vect{x}_{\mathcal{E}}^{c\,i\,d}
    -
    \mat{A}\vect{y}
  \right)^{\top}
  \mat{W}^{c\,i\,d}
  \left(
    \vect{x}_{\mathcal{E}}^{c\,i\,d}
    -
    \mat{A}\vect{y}
  \right),
  \qquad
  \mat{G}_{\star}^{\top}\vect{y}=\vect{0}.
\]
For a connected usable graph, the constraint is simply
\[
  \vect{w}_{\star}^{\top}\vect{y}=0.
\]

This recovered node coordinate becomes the node feature
\[
  s_u^{c\,i\,d}
  =
  \widehat{y}_{u}^{c\,i\,d},
  \qquad
  d\in\mathcal{J}_{\mathrm{price}}.
\]

Let
\[
  \mathcal{C}_{\mathrm{rec}}^{c\,i\,d}(t)
\]
be the set of recoverable connected components for price coordinate $d$.  Define
the price-coordinate node mask by
\[
  {}_{s}m_{\mathrm{price}}^{c\,i\,u\,d}(t)
  =
  \mathbf{1}
  \left\{
    u
    \in
    \bigcup_{a\in\mathcal{C}_{\mathrm{rec}}^{c\,i\,d}(t)}
    \mathcal{V}_{(a)}^{c\,i\,d}(t)
  \right\}.
\]
When this mask is zero, the corresponding recovered price coordinate
$s_u^{c\,i\,d}(t)$ is not treated as observed node evidence.

Define the graph precision
\[
  \mat{H}^{c\,i\,d}
  =
  \mat{A}^{\top}
  \mat{W}^{c\,i\,d}
  \mat{A}
  \in
  \R^{N\times N}.
\]
Then the constrained solve may be written as the KKT system
\[
  \begin{bmatrix}
    \mat{H}^{c\,i\,d} & \mat{G}_{\star}\\
    \mat{G}_{\star}^{\top} & \mat{0}
  \end{bmatrix}
  \begin{bmatrix}
    \widehat{\vect{y}}^{c\,i\,d}\\
    \vect{\lambda}_{\star}^{c\,i\,d}
  \end{bmatrix}
  =
  \begin{bmatrix}
    \mat{A}^{\top}
    \mat{W}^{c\,i\,d}
    \vect{x}_{\mathcal{E}}^{c\,i\,d}\\
    \vect{0}
  \end{bmatrix}.
\]
The vector $\vect{\lambda}_{\star}^{c\,i\,d}$ contains the Lagrange multipliers
for the component reference constraints.  The node feature is
$\widehat{\vect{y}}^{c\,i\,d}$.

When the edge precision is diagonal,
\[
  \mat{W}^{c\,i\,d}
  =
  \mat{\Gamma}^{c\,i\,d},
\]
the matrix
\[
  \mat{L}_{\gamma}^{c\,i\,d}
  :=
  \mat{A}^{\top}
  \mat{\Gamma}^{c\,i\,d}
  \mat{A}
  \in
  \R^{N\times N}
\]
is the confidence-weighted graph Laplacian of the usable graph for that
coordinate.  In this case,
\[
  \mat{H}^{c\,i\,d}
  =
  \mat{L}_{\gamma}^{c\,i\,d}.
\]

When a full edge covariance is used,
\[
  \mat{W}^{c\,i\,d}
  =
  \left(
    \mat{C}_{x}^{c\,i\,d}
  \right)^{-1},
\]
the same matrix
\[
  \mat{H}^{c\,i\,d}
  =
  \mat{A}^{\top}
  \mat{W}^{c\,i\,d}
  \mat{A}
\]
is the covariance-aware graph precision.  It plays the same role as the
Laplacian in the recovery solve, but it may include cross-edge uncertainty
rather than only independent edge weights.

Both the Laplacian and the covariance-aware graph precision have the same
structural ambiguity: on each connected component, a constant node vector is
invisible to edge differences.  The reference constraints in
$\mat{G}_{\star}^{\top}\vect{y}=\vect{0}$ remove those additive null directions.

The close coordinate is the cleanest scalar return coordinate.  The open, high,
and low coordinates are also price-ratio features, so the same lift is useful
for them.  However, high and low are candle extrema, so their node lift should
be read as a feature-wise reconstruction rather than a claim that all edge
extrema were synchronized through one instantaneous node value.

\subsection{Price Residual Sidecar}
\label{subsec:price-residual-sidecar}

For every lifted price coordinate, keep the edge residual
\[
  \vect{\eps}_{\mathrm{price}}^{c\,i\,d}
  =
  \vect{x}_{\mathcal{E}}^{c\,i\,d}
  -
  \mat{A}
  \widehat{\vect{y}}^{c\,i\,d}
  \in
  \R^{L_{\mathrm{use}}^{c\,i\,d}},
  \qquad
  d\in\mathcal{J}_{\mathrm{price}}.
\]
The weighted residual energy is
\[
  \mathcal{Q}_{\mathrm{price}}^{c\,i\,d}
  =
  \left(
    \vect{\eps}_{\mathrm{price}}^{c\,i\,d}
  \right)^{\top}
  \mat{W}^{c\,i\,d}
  \vect{\eps}_{\mathrm{price}}^{c\,i\,d}
  \in
  \Rnonneg.
\]

This residual is not a nuisance.  It is the edge-local part that cannot be
explained by one global node-potential vector.  On cycles, it records
non-integrability or cycle tension.  It may later be useful for arbitrage-like
disagreement diagnostics, but executable trading decisions remain outside this
notation layer.

\calloutbox{Residuals are preserved.}{\\%
  Node lifting does not throw away the edge world.  It decomposes price-like
  edge fields into an integrable node-potential part and a residual edge part:
  \[
    \vect{x}_{\mathcal{E}}^{c\,i\,d}
    =
    \mat{A}\widehat{\vect{y}}^{c\,i\,d}
    +
    \vect{\eps}_{\mathrm{price}}^{c\,i\,d}.
  \]
  The node tensor carries the semantic global state.  The residual sidecar
  carries the remaining edge tension.
}

The residual sidecar also receives a mask.  For $e=(u,v)$ and
$d\in\mathcal{J}_{\mathrm{price}}$, define
\[
  {}_{\eps}m_{\mathrm{price},e}^{c\,i\,d}(t)
  =
  {}_{x}m_{e}^{c\,i\,d}(t)
  \cdot
  {}_{s}m_{\mathrm{price}}^{c\,i\,u\,d}(t)
  \cdot
  {}_{s}m_{\mathrm{price}}^{c\,i\,v\,d}(t).
\]
This says that a price residual is valid only when the original edge coordinate
was valid and both endpoint nodes belong to a recoverable price component.
Invalid residuals are not used by downstream residual diagnostics or later
trading-engine logic.

\subsection{Activity-Coordinate Lift}
\label{subsec:activity-coordinate-lift}

Activity coordinates are not differences of node potentials.  They are
nonnegative quantities attached to a traded relation.  Therefore they are lifted
by endpoint aggregation.

A direct sum of incident edge activity is meaningful, but it is also sensitive
to the number of usable edges touching a node.  Since the same representation
engine is intended to operate across all nodes, NodeLift separates total node
activity from support-normalized node activity intensity.

For
\[
  d\in\mathcal{J}_{\mathrm{act}},
\]
the normalized coordinate has the form
\[
  x_{e}^{c\,i\,d}
  =
  \log
  \left(
    1+a_{e}^{c\,i\,d}
  \right).
\]
For observed normalized input data, recover the raw nonnegative activity by
\[
  a_{e}^{c\,i\,d}
  =
  \exp
  \left(
    x_{e}^{c\,i\,d}
  \right)
  -
  1
  \in
  \Rnonneg.
\]
The inverse activity map is applied only to valid normalized activity
coordinates.

For each activity coordinate, define the endpoint assignment matrix
\[
  \mat{M}^{(d)}
  =
  \begin{cases}
    \mat{B}_{\mathcal{E}}, & d=5 \quad \text{base volume},\\
    \mat{Q}_{\mathcal{E}}, & d=6 \quad \text{quote volume},\\
    \mat{U}_{\mathcal{E}}, & d=7 \quad \text{incident trades},\\
    \mat{B}_{\mathcal{E}}, & d=8 \quad \text{taker buy base volume},\\
    \mat{Q}_{\mathcal{E}}, & d=9 \quad \text{taker buy quote volume}.
  \end{cases}
\]
The valid support count at node $u$ is
\[
  \kappa_{u}^{c\,i\,d}(t)
  =
  \sum_{e\in\mathcal{E}}
  M_{e}^{(d),u}
  \cdot
  {}_{x}m_{e}^{c\,i\,d}(t)
  \in
  \N_{\ge 0}.
\]
This counts how many valid assigned edge coordinates contribute to node $u$ for
activity coordinate $d$.

The summed raw node activity is
\[
  \nu_{\Sigma,u}^{c\,i\,d}(t)
  =
  \sum_{e\in\mathcal{E}}
  M_{e}^{(d),u}
  \cdot
  {}_{x}m_{e}^{c\,i\,d}(t)
  \cdot
  a_{e}^{c\,i\,d}(t)
  \in
  \Rnonneg.
\]
This total is useful as a liquidity or execution-side quantity, but it is
degree-sensitive.

The support-normalized activity intensity is
\[
  \bar{\nu}_{u}^{c\,i\,d}(t)
  =
  \frac{
    \nu_{\Sigma,u}^{c\,i\,d}(t)
  }{
    \max
    \left(
      1,
      \kappa_{u}^{c\,i\,d}(t)
    \right)
  }
  \in
  \Rnonneg.
\]
This quantity measures average activity per valid assigned edge.  It is less
sensitive to the number of markets touching the node and is therefore the
default activity value exported to the shared node representation engine.

The normalized node activity feature is
\[
  s_u^{c\,i\,d}(t)
  =
  \log
  \left(
    1+\bar{\nu}_{u}^{c\,i\,d}(t)
  \right),
  \qquad
  d\in\mathcal{J}_{\mathrm{act}}.
\]
The activity-coordinate node mask is
\[
  {}_{s}m_{\mathrm{act}}^{c\,i\,u\,d}(t)
  =
  \mathbf{1}
  \left\{
    \kappa_{u}^{c\,i\,d}(t)>0
  \right\},
  \qquad
  d\in\mathcal{J}_{\mathrm{act}}.
\]
This mask says that node $u$ has valid lifted activity coordinate $d$ whenever
at least one usable assigned edge contributes to that node.

A useful optional diagnostic is the activity coverage fraction
\[
  \rho_{u}^{c\,i\,d}(t)
  =
  \frac{
    \kappa_{u}^{c\,i\,d}(t)
  }{
    \max
    \left(
      1,
      \sum_{e\in\mathcal{E}}M_{e}^{(d),u}
    \right)
  }
  \in[0,1].
\]
The mask records whether any usable activity evidence exists.  The coverage
fraction records how much of the staged incident activity support was actually
available.

The trade-count lift above uses the full incident count at each endpoint.  If a
future diagnostic needs graph-wide conservation of total trade counts, one may
replace $\mat{U}_{\mathcal{E}}^{\top}$ by
$\frac{1}{2}\mat{U}_{\mathcal{E}}^{\top}$.  The full incident version is the
default here because the node feature asks: ``how much activity touched this
node?''

\calloutbox{Activity normalization.}{\\%
  The summed activity $\nu_{\Sigma,u}^{c\,i\,d}$ answers: ``how much total
  activity touched this node?''  The normalized intensity
  $\bar{\nu}_{u}^{c\,i\,d}$ answers: ``how much activity occurred per valid
  assigned edge?''  The shared node representation uses the normalized
  intensity by default, while the total activity and coverage may be kept as
  sidecars.
}

\subsection{The Full Node Feature Tensor and Exported Masks}
\label{subsec:full-node-feature-tensor-and-exported-masks}

Combining the price lift and activity lift gives the full node feature matrix
\[
  \tensor{s}^{c\,i}(t)
  \in
  \R^{N\times D_y}.
\]
For the full normalized kline feature lift,
\[
  D_y=D_x=9.
\]
Coordinatewise,
\[
  s_u^{c\,i\,d}(t)
  =
  \begin{cases}
    \widehat{y}_{u}^{c\,i\,d}(t),
    & d\in\mathcal{J}_{\mathrm{price}},\\[0.8em]
    \log
    \left(
      1+\bar{\nu}_{u}^{c\,i\,d}(t)
    \right),
    & d\in\mathcal{J}_{\mathrm{act}}.
  \end{cases}
\]

The full coordinatewise node mask is
\[
  {}_{s}m^{c\,i\,u\,d}(t)
  =
  \begin{cases}
    {}_{s}m_{\mathrm{price}}^{c\,i\,u\,d}(t),
    & d\in\mathcal{J}_{\mathrm{price}},\\[0.6em]
    {}_{s}m_{\mathrm{act}}^{c\,i\,u\,d}(t),
    & d\in\mathcal{J}_{\mathrm{act}}.
  \end{cases}
\]
When a downstream node encoder accepts only a node-level mask, derive a coarse
mask from the coordinatewise one.  An any-valid mask is
\[
  {}_{s}m_{\mathrm{any}}^{c\,i\,u}(t)
  =
  \mathbf{1}
  \left\{
    \sum_{d=1}^{D_y}
    {}_{s}m^{c\,i\,u\,d}(t)
    >
    0
  \right\}.
\]
An all-required mask for a chosen coordinate set
$\mathcal{J}_{\mathrm{req}}\subseteq\{1,\ldots,D_y\}$ is
\[
  {}_{s}m_{\mathrm{all}}^{c\,i\,u}(t)
  =
  \prod_{d\in\mathcal{J}_{\mathrm{req}}}
  {}_{s}m^{c\,i\,u\,d}(t).
\]
The implementation should name which coarse rule is used.  If the node encoder
does not accept coordinatewise masks, the all-required rule is safer because it
does not treat missing feature coordinates as observed values.

Over a full observed input window, the node-lifted input tensor is
\[
  \tensor{s}_{x}(t)
  \in
  \R^{C\times H_x\times N\times D_y}.
\]
The corresponding coordinatewise node mask is
\[
  {}_{s}\tensor{m}_{x}(t)
  \in
  \{0,1\}^{C\times H_x\times N\times D_y}.
\]
A coarser node mask may also be exported when needed:
\[
  {}_{s}\tensor{m}_{x,\mathrm{any}}(t)
  \in
  \{0,1\}^{C\times H_x\times N},
\]
or
\[
  {}_{s}\tensor{m}_{x,\mathrm{all}}(t)
  \in
  \{0,1\}^{C\times H_x\times N}.
\]
This is the NodeLift output used by the node-first representation path.

\subsection{NodeLift as a Procedure}
\label{subsec:nodelift-procedure}

For one anchor time $t$, one channel $c$, and one input position $i$, the
concrete procedure is:

\begin{enumerate}
  \item Start from the observed edge feature slices and input masks
  \[
    \left\{
      \tensor{x}_{e}^{c\,i}(t),
      {}_{x}m_{e}^{c\,i}(t)
      :
      e\in\mathcal{E}
    \right\}.
  \]

  \item Build the coordinatewise edge masks
  \[
    {}_{x}m_{e}^{c\,i\,d}(t)
    =
    {}_{x}m_{e}^{c\,i}(t)
    \cdot
    \mathbf{1}
    \left\{
      x_{e}^{c\,i\,d}(t)
      \text{ is finite}
    \right\}.
  \]
  These masks determine the coordinatewise usable edge sets
  \[
    \mathcal{E}_{\mathrm{use}}^{c\,i\,d}(t)
    =
    \left\{
      e\in\mathcal{E}:
      {}_{x}m_{e}^{c\,i\,d}(t)=1
    \right\}.
  \]

  \item Build the full active graph matrices
  \[
    \mat{A}_{\mathcal{E}},
    \qquad
    \mat{B}_{\mathcal{E}},
    \qquad
    \mat{Q}_{\mathcal{E}},
    \qquad
    \mat{U}_{\mathcal{E}}
    =
    \mat{B}_{\mathcal{E}}
    +
    \mat{Q}_{\mathcal{E}}.
  \]
  For each coordinate, restrict these matrices to the rows selected by
  $\mathcal{E}_{\mathrm{use}}^{c\,i\,d}(t)$.

  \item For each price coordinate
  $d\in\mathcal{J}_{\mathrm{price}}$, find the connected components of the
  underlying undirected usable graph induced by
  $\mathcal{E}_{\mathrm{use}}^{c\,i\,d}(t)$.  For each recoverable component,
  choose or update a synthetic reference
  $\vect{w}_{\star}^{(a)}$, and assemble the constraint matrix
  \[
    \mat{G}_{\star}^{c\,i\,d}(t).
  \]

  \item For each price coordinate
  $d\in\mathcal{J}_{\mathrm{price}}$, solve the gauge-fixed weighted
  least-squares problem on the recoverable components:
  \[
    \vect{x}_{\mathcal{E}}^{c\,i\,d}
    \approx
    \mat{A}\vect{y}^{c\,i\,d},
    \qquad
    \left(
      \mat{G}_{\star}^{c\,i\,d}
    \right)^{\top}
    \vect{y}^{c\,i\,d}
    =
    \vect{0}.
  \]
  Store the recovered node coordinate
  \[
    \widehat{\vect{y}}^{c\,i\,d}
  \]
  and the price residual
  \[
    \vect{\eps}_{\mathrm{price}}^{c\,i\,d}.
  \]
  Also store the price-coordinate node mask
  \[
    {}_{s}m_{\mathrm{price}}^{c\,i\,u\,d}(t)
  \]
  and the residual mask
  \[
    {}_{\eps}m_{\mathrm{price},e}^{c\,i\,d}(t).
  \]

	  \item For each activity coordinate
	  $d\in\mathcal{J}_{\mathrm{act}}$, undo the $\log(1+\cdot)$ transform on valid
	  edge coordinates:
  \[
    a_{e}^{c\,i\,d}(t)
    =
    \exp
    \left(
      x_{e}^{c\,i\,d}(t)
    \right)
	    -
	    1.
	  \]
	  Aggregate valid raw activity to endpoints:
	  \[
	    \nu_{\Sigma,u}^{c\,i\,d}(t)
	    =
	    \sum_{e\in\mathcal{E}}
	    M_{e}^{(d),u}
	    \cdot
	    {}_{x}m_{e}^{c\,i\,d}(t)
	    \cdot
	    a_{e}^{c\,i\,d}(t).
	  \]
	  Compute the valid support count:
	  \[
	    \kappa_{u}^{c\,i\,d}(t)
	    =
	    \sum_{e\in\mathcal{E}}
	    M_{e}^{(d),u}
	    \cdot
	    {}_{x}m_{e}^{c\,i\,d}(t).
	  \]
	  Compute the support-normalized activity intensity:
	  \[
	    \bar{\nu}_{u}^{c\,i\,d}(t)
	    =
	    \frac{
	      \nu_{\Sigma,u}^{c\,i\,d}(t)
	    }{
	      \max
	      \left(
	        1,
	        \kappa_{u}^{c\,i\,d}(t)
	      \right)
	    }.
	  \]
	  Export the normalized node activity feature
	  \[
	    s_u^{c\,i\,d}(t)
	    =
	    \log
	    \left(
	      1+\bar{\nu}_{u}^{c\,i\,d}(t)
	    \right).
	  \]
	  Store the activity-coordinate node mask
	  \[
	    {}_{s}m_{\mathrm{act}}^{c\,i\,u\,d}(t)
    =
    \mathbf{1}
	    \left\{
	      \kappa_{u}^{c\,i\,d}(t)>0
	    \right\}.
	  \]
	  Optionally store the summed activity
	  $\nu_{\Sigma,u}^{c\,i\,d}(t)$ and coverage
	  $\rho_u^{c\,i\,d}(t)$ as sidecars.

  \item Assemble the node feature matrix
  \[
    \tensor{s}^{c\,i}(t)
    \in
    \R^{N\times D_y}.
  \]
  Coordinatewise,
  \[
    s_u^{c\,i\,d}(t)
    =
    \begin{cases}
      \widehat{y}_{u}^{c\,i\,d}(t),
      & d\in\mathcal{J}_{\mathrm{price}},\\[0.8em]
	      \log
	      \left(
	        1+\bar{\nu}_{u}^{c\,i\,d}(t)
	      \right),
      & d\in\mathcal{J}_{\mathrm{act}}.
    \end{cases}
  \]

  \item Assemble the coordinatewise node mask
  \[
    {}_{s}\tensor{m}^{c\,i}(t)
    \in
    \{0,1\}^{N\times D_y},
  \]
  where
  \[
    {}_{s}m^{c\,i\,u\,d}(t)
    =
    \begin{cases}
      {}_{s}m_{\mathrm{price}}^{c\,i\,u\,d}(t),
      & d\in\mathcal{J}_{\mathrm{price}},\\[0.6em]
      {}_{s}m_{\mathrm{act}}^{c\,i\,u\,d}(t),
      & d\in\mathcal{J}_{\mathrm{act}}.
    \end{cases}
  \]

  \item Optionally derive a coarse node mask
  \[
    {}_{s}m_{\mathrm{any}}^{c\,i\,u}(t)
    \quad\text{or}\quad
    {}_{s}m_{\mathrm{all}}^{c\,i\,u}(t),
  \]
  depending on what the downstream node encoder accepts.

  \item Export the node feature matrix, node masks, and residual sidecar:
  \[
    \left(
      \tensor{s}^{c\,i}(t),
      {}_{s}\tensor{m}^{c\,i}(t),
      \tensor{\eps}_{\mathrm{price}}^{c\,i}(t),
      {}_{\eps}\tensor{m}_{\mathrm{price}}^{c\,i}(t)
    \right).
  \]
\end{enumerate}

In compact form:
\[
  \operatorname{NodeLift}
  \left[
    \left\{
      \tensor{x}_{e}^{c\,i}(t),
      {}_{x}m_{e}^{c\,i}(t)
      :
      e\in\mathcal{E}
    \right\},
    \mat{A}_{\mathcal{E}},
    \mat{B}_{\mathcal{E}},
    \mat{Q}_{\mathcal{E}},
    \mat{G}_{\star}
  \right]
  =
  \left(
    \tensor{s}^{c\,i}(t),
    {}_{s}\tensor{m}^{c\,i}(t),
    \tensor{\eps}_{\mathrm{price}}^{c\,i}(t),
    {}_{\eps}\tensor{m}_{\mathrm{price}}^{c\,i}(t)
  \right).
\]

This is the edge-to-node lift used by the node-first representation path.

\subsection{Node Geometry After Lifting}
\label{subsec:node-geometry-after-lifting}

For price coordinates, the recovery solve also gives a local geometry on node
states.  For one coordinate $d\in\mathcal{J}_{\mathrm{price}}$, define
\[
  \mat{H}^{c\,i\,d}
  =
  \mat{A}^{\top}
  \mat{W}^{c\,i\,d}
  \mat{A}
  \in
  \R^{N\times N}.
\]
Let
\[
  \mat{B}_{\star}
\]
have columns spanning the synthetic-gauge free subspace:
\[
  \mat{G}_{\star}^{\top}\mat{B}_{\star}=\mat{0}.
\]
If the usable graph is connected, then
\[
  \mat{B}_{\star}\in\R^{N\times(N-1)}.
\]
If the usable graph has $C_{\mathrm{comp}}$ connected components, then the free
dimension is
\[
  N_{\mathrm{use}}-C_{\mathrm{comp}}
\]
on the touched usable nodes.

The free-coordinate graph precision is
\[
  \mat{H}_{\star}^{c\,i\,d}
  =
  \mat{B}_{\star}^{\top}
  \mat{H}^{c\,i\,d}
  \mat{B}_{\star}.
\]
When well posed, the corresponding node covariance in free coordinates is
\[
  \mat{C}_{\theta}^{c\,i\,d}
  \approx
  \left(
    \mat{H}_{\star}^{c\,i\,d}
  \right)^{-1}.
\]
Lifting back to node coordinates gives
\[
  \mat{C}_{y}^{c\,i\,d}
  =
  \mat{B}_{\star}
  \mat{C}_{\theta}^{c\,i\,d}
  \mat{B}_{\star}^{\top}.
\]
This covariance is understood on the synthetic-gauge subspace.  The removed
constant directions are not economic uncertainty modes; they are only the
additive symmetries of the edge-difference equations.

For two node states $\vect{y},\vect{y}'$ in the same gauge, the Mahalanobis
distance is
\[
  d_{\star}(\vect{y},\vect{y}')^2
  =
  (\vect{y}-\vect{y}')^{\top}
  \left(
    \mat{C}_{y}^{c\,i\,d}
  \right)^{+}
  (\vect{y}-\vect{y}')
  \in
  \Rnonneg.
\]
This measures distance in units of node uncertainty rather than raw coordinate
scale.

\subsection{Modes of Node Motion}
\label{subsec:modes-of-node-motion}

When the covariance is available in free synthetic-gauge coordinates, its
eigenvectors describe collective modes of node motion:
\[
  \mat{C}_{\theta}^{c\,i\,d}
  =
  \mat{V}_{\star}^{c\,i\,d}
  \Lambda_{\star}^{c\,i\,d}
  \left(
    \mat{V}_{\star}^{c\,i\,d}
  \right)^{\top}.
\]
The columns of $\mat{V}_{\star}^{c\,i\,d}$ are principal directions in the free
node coordinates, and the diagonal entries of $\Lambda_{\star}^{c\,i\,d}$
measure variance along those directions.

The corresponding node-coordinate modes are
\[
  \mat{B}_{\star}
  \mat{V}_{\star}^{c\,i\,d}.
\]
These modes are optional readouts of the recovered node geometry.  They are not
required for the basic node lift, but they become useful when the next semantic
layer wants lower-dimensional collective node factors.

\subsection{What This Chapter Gives to Cuwacunu}
\label{subsec:what-nodelift-gives-cuwacunu}

The chapter defines a graph transformation:
\[
  \text{edge feature tensor}
  \quad
  \longrightarrow
  \quad
  \text{node feature tensor}
  \quad
  +
  \quad
  \text{price residual sidecar}.
\]

The transformation is not a neural network.  It is a deterministic graph lift
defined by the market graph, the feature registry, the masks, the edge precision
weights, and the synthetic reference gauge.

This gives Cuwacunu the input-side node-first path:
\[
  \tensor{x}_{u\,v}(t)
  \longrightarrow
  \tensor{s}_{x}(t)
  \longrightarrow
  \text{node representation learning}.
\]

\calloutbox{Modeling route enabled by NodeLift.}{\\%
  In the node-first route, observed edge fields are first lifted to node tensors
  and a later model learns node representations directly.  The hybrid variant
  keeps node tensors and residual edge sidecars together, preserving both
  semantic node state and edge-local cycle tension.
}

The node-first path does not need to discard edge information.  It keeps
\[
  \tensor{\eps}_{\mathrm{price}}(t)
\]
as a sidecar for the residual edge tension.  Later systems may also keep edge
activity sidecars when edge-local liquidity, routing, or execution detail
matters.

\calloutbox{End of the NodeLift layer.}{\\%
  At this point the document has defined a concrete procedure for lifting
  normalized edge-carried market features into node-carried market features.
  Price coordinates are recovered through incidence-based balance equations.
  In the independent-confidence case, the recovery matrix is the
  confidence-weighted graph Laplacian.  Activity coordinates are recovered
  through endpoint aggregation.  Synthetic gauges make node coordinates
  interpretable, masks are lifted together with values, and residuals preserve
  the edge-local tension that the node state cannot explain.  Trading policy,
  execution rules, and portfolio decision logic remain outside this notation
  layer.
}
